% STATUS: written 2026-11-05 (Kadison 1966 and Sakai 1966 read in full;
% excerpts in excerpts/Kadison1966.txt and excerpts/Sakai1966.txt).
% Additions of 2026-11-06 from Ayupov--Arzikulov 2011 (pp. 1--7),
% Ayupov--Kudaybergenov 2014 (pp. 1--2, 11), Hadwin--Li--Li--Ma 2013 (pp. 1--3),
% Albeverio--Ayupov--Kudaybergenov 2007 (pp. 1, 5, 11) and
% Albeverio--Ayupov--Kudaybergenov--Kalandarov 2008 (pp. 1, 5, 10).
% Additions of 2026-11-07 from Kadison--Liu 2014 (pp. 2087--2091) and
% Yu--Zhang 2012 (pp. 1979--1980, 1990--1991).
% Sources: Kadison1966, Sakai1966, Sakai1960Kaplansky, Liu2025TwoLocal,
% AyupovArzikulov2011TwoLocal, AyupovKudaybergenov2014Matrix,
% HadwinLiLiMa2013, AlbeverioAyupovKudaybergenov2007,
% AlbeverioAyupovKudaybergenovKalandarov2008, HaoChen2026Local,
% KadisonLiu2014Murray, YuZhang2012StarLie.
\section{Kadison--Sakai: derivations of von Neumann algebras}\label{sec:banach-operator:kadison-sakai}

Section~\ref{sec:banach-operator:singer-wermer} ended with the commutative
$C^*$-algebras: a derivation of a commutative $C^*$-algebra is zero
\cite[p.~280]{Kadison1966}. On that basis Kaplansky proved that every
derivation of a type I von Neumann algebra is inner
\cite[p.~280]{Kadison1966}. He proposed
that every derivation of a $C^*$-algebra is continuous
\cite[p.~280]{Kadison1966}. Sakai proved the continuity in 1960
\cite[p.~31]{Sakai1960Kaplansky}. Miles then noted that each derivation of
a $C^*$-algebra is induced by an operator in the weak closure of some
faithful representation \cite[pp.~273--274]{Sakai1966}. The remaining
question was whether a von Neumann algebra can have an outer derivation at
all \cite[p.~273]{Sakai1966}. The vocabulary is that of Sakai
\cite[pp.~273--274]{Sakai1966}.

\begin{definition}[{\cite[pp.~273--274]{Sakai1966}}]\label{def:banach:inner}
Let $A$ be a $C^*$-algebra and $D$ a derivation of $A$. The derivation $D$
is an \emph{inner derivation} if there is an element $a$ in $A$ such that
$D(x)=[a,x]=ax-xa$ for $x$ in $A$. Otherwise it is an \emph{outer
derivation}.
\end{definition}

A derivation of $A$ is \emph{spatial} when a bounded operator $B$ on $H$
implements it, so that $\delta(A)=BA-AB$ for each $A$ in
$A$~\cite[p.~281]{Kadison1966}. Kadison began with the centre~\cite[Theorem~2, p.~281]{Kadison1966}. The first step
controls how a derivation behaves
 there~ (Theorem~\ref{thm:banach:center})\cite[Theorem~2, p.~281]{Kadison1966}.

\begin{theorem}[{\cite[Theorem~2, p.~281]{Kadison1966}}]\label{thm:banach:center}
Each derivation of a $C^*$-algebra annihilates its center.
\end{theorem}

Two facts drive the proof~\cite[p.~281]{Kadison1966}. A pure state is
multiplicative on the centre, and a derivation kills every element on which a
definite state vanishes to second order~\cite[p.~281]{Kadison1966}. With the centre out of the way,
extension to the weak closure gives the spatial
theorem~\cite[p.~281]{Kadison1966}.

\begin{theorem}[{\cite[Theorem~4, p.~283]{Kadison1966}}]\label{thm:banach:spatial}
Each derivation $\delta$ of a $C^*$-algebra $A$ acting on the Hilbert space
$H$ is spatial. We may choose $B$ commuting with an assigned maximal
abelian subalgebra of $A'$ or with an assigned mixing von Neumann
subalgebra of $A'$ so that $\operatorname{ad} B\,\vert\,A=\delta$.
\end{theorem}

For hyperfinite von Neumann algebras the implementing operator already lies
in the algebra (Theorem~\ref{thm:banach:hyperfinite}).

\begin{theorem}[{\cite[Theorem~7, p.~285]{Kadison1966}}]\label{thm:banach:hyperfinite}
Each derivation of a mixing von Neumann algebra is inner. In particular,
each derivation of a hyperfinite von Neumann algebra is inner.
\end{theorem}

Sakai then closed the question~\cite[p.~274]{Sakai1966}. The argument takes a least element of a
convex set in the double commutant~\cite[p.~290]{Kadison1966}. The paper of
Kadison presents that step as the completion of its own
program~\cite[p.~290]{Kadison1966}.

\begin{theorem}[{\cite[Theorem~1, p.~274]{Sakai1966}}]\label{thm:banach:sakaiinner}
Every derivation of a $W^*$-algebra is inner.
\end{theorem}

Theorem~\ref{thm:banach:sakaiinner} answers the conjecture of Kadison in the
affirmative~\cite[p.~274]{Sakai1966}. Kadison states the same result in the
language of von Neumann algebras, and attributes it to Sakai, who completed
the arguments of his paper~\cite[Theorem~15, p.~290]{Kadison1966}.

\begin{theorem}[{\cite[Theorem~15, p.~290]{Kadison1966}}]\label{thm:banach:kadison-inner}
Each derivation of a von Neumann algebra is inner.
\end{theorem}

A von Neumann algebra is a $W^*$-algebra~\cite[p.~2]{Sakai1991Dynamical}, so
the statement of Theorem~\ref{thm:banach:kadison-inner} and the statement of
Theorem~\ref{thm:banach:sakaiinner} agree. Kadison also records the
consequences for quantum physics
\cite[p.~291]{Kadison1966}. Take a strongly continuous one-parameter unitary
group whose spectrum is bounded below. If it induces automorphisms of a von
Neumann algebra, the automorphisms are inner \cite[p.~291]{Kadison1966}. In
physical language, energy and momentum are observable without the assumption
of a vacuum state \cite[p.~291]{Kadison1966}.

The $C^*$-algebra form places the
implementing element in the weak closure
\cite[Theorem~2, p.~274]{Sakai1966}.

\begin{theorem}[{\cite[Theorem~2, p.~274]{Sakai1966}}]\label{thm:banach:sakaiclosure}
Let $A$ be a $C^*$-algebra on a Hilbert space $H$, $D$ a derivation on $A$,
$A^-$ the weak closure of $A$ on $H$. Then there is a bounded operator $b$
belonging to $A^-$ such that $D(x)=[b,x]$ for all $x$ in $A$.
\end{theorem}

The theorems of this section lean on the automatic continuity of Sakai from
1960 \cite[p.~283]{Kadison1966}. The proof of Kadison for Theorem~\ref{thm:banach:spatial} opens by
citing it for boundedness \cite[p.~283]{Kadison1966}. Continuity has a story of its own,
told in Section~\ref{sec:banach-operator:automatic-continuity}.

A recent line of work drops the assumption that the map is linear at
all~\cite[p.~909]{Liu2025TwoLocal}. The
definition is due to Šemrl and is recalled by Liu in the form that
follows~\cite[Definition~1.2, p.~909]{Liu2025TwoLocal}.

\begin{definition}[{\cite[Definition~1.2, p.~909]{Liu2025TwoLocal}}]\label{def:banach:two-local}
Let $A$ be an algebra. A map $\delta\colon A \to A$, not assumed linear, is a
\emph{2-local derivation} when for all $A, B \in A$ there is a linear
derivation $\delta_{A,B}\colon A \to A$, depending on the two elements, with
\[
  \delta(A)\,B + A\,\delta(B) = \delta_{A,B}(AB).
\]
\end{definition}

This definition generalizes the earlier one, in which only
$\delta(A) = \delta_{A,B}(A)$ and $\delta(B) = \delta_{A,B}(B)$ are asked for
\cite[p.~909]{Liu2025TwoLocal}. Molnár proved that every 2-local derivation of
a standard operator algebra into the bounded operators is a linear
derivation~\cite[p.~910]{Liu2025TwoLocal}. He then asked whether the same
holds on an arbitrary von Neumann
algebra~\cite[p.~910]{Liu2025TwoLocal}. Liu answers the question on the
semi-finite factors~ (Theorem~\ref{thm:banach:two-local})\cite[Theorem~2.1, p.~910]{Liu2025TwoLocal}.

\begin{theorem}[{\cite[Theorem~2.1, p.~910]{Liu2025TwoLocal}}]\label{thm:banach:two-local}
Let $A$ be a semi-finite factor von Neumann algebra, and let
$\delta\colon A \to A$ be a map such that for all $A, B \in A$ there is a
linear derivation $\delta_{A,B}\colon A \to A$ with
$\delta(A)B + A\delta(B) = \delta_{A,B}(AB)$. Then $\delta$ is a linear
derivation.
\end{theorem}

The proof uses the trace~\cite[p.~910]{Liu2025TwoLocal}. Every linear
derivation of such an algebra is inner~\cite[p.~910]{Liu2025TwoLocal}, which
is the theorem of Sakai stated below as
Theorem~\ref{thm:banach:sakaiclosure}. Take a trace-class element $M$. The
inner form makes $\delta(A)M + A\delta(M)$ have trace
zero~\cite[p.~910]{Liu2025TwoLocal}. Linearity then follows because the trace
is faithful~\cite[p.~910]{Liu2025TwoLocal}.

On von Neumann algebras the answer is now complete. Ayupov and Arzikulov
work with the original notion of \v{S}emrl, where the two values alone are
matched \cite[p.~1]{AyupovArzikulov2011TwoLocal}. For a von Neumann algebra
every derivation is inner, so the condition reads as follows
\cite[p.~2]{AyupovArzikulov2011TwoLocal}. For each pair $x,y$ of elements
there is an element $a$ of the algebra with $\Delta(x)=ax-xa$ and
$\Delta(y)=ay-ya$ \cite[p.~2]{AyupovArzikulov2011TwoLocal}.

\begin{theorem}[{\cite[Theorem~5, p.~7]{AyupovArzikulov2011TwoLocal}}]\label{thm:banach:two-local-typeI}
Let $M$ be a von Neumann algebra of type I and let $\Delta\colon M\to M$ be
a 2-local derivation. Then $\Delta$ is a derivation.
\end{theorem}

The proof splits the algebra into its homogeneous summands, each of the form
$C(X)\otimes M_n(\mathbb{C})$ \cite[p.~7]{AyupovArzikulov2011TwoLocal}. On
one summand the additivity follows from the matrix units
\cite[Theorem~1, p.~3]{AyupovArzikulov2011TwoLocal}. A Jordan derivation on
a semiprime algebra is a derivation, so additivity is enough
\cite[p.~3]{AyupovArzikulov2011TwoLocal}.

The same authors carry the result to matrix algebras over a semiprime Banach
algebra \cite[p.~1]{AyupovKudaybergenov2014Matrix}. An algebra has the
inner derivation property when every derivation of it is inner (Theorem~\ref{thm:banach:two-local-matrix})
\cite[p.~2]{AyupovKudaybergenov2014Matrix}.\begin{theorem}[{\cite[Theorem~2.1, p.~2]{AyupovKudaybergenov2014Matrix}}]\label{thm:banach:two-local-matrix}
Let $A$ be a unital semiprime Banach algebra with the inner derivation
property and let $M_{2^n}(A)$ be the algebra of $2^n\times 2^n$ matrices
over $A$, with $n\geq 2$. Then every 2-local derivation of $M_{2^n}(A)$ is
a derivation.
\end{theorem}

The matrix case passes to $\mathrm{AW}^*$-algebras
\cite[p.~11]{AyupovKudaybergenov2014Matrix}. Every
$\mathrm{AW}^*$-algebra has a central decomposition into summands that are
matrix algebras over abelian ones, and the theorem applies to each summand
\cite[p.~11]{AyupovKudaybergenov2014Matrix}. So every 2-local derivation
of an $\mathrm{AW}^*$-algebra is a derivation
\cite[Theorem~3.1, p.~11]{AyupovKudaybergenov2014Matrix}. Together with
Theorem~\ref{thm:banach:two-local-typeI} this settles the question of
Moln\'ar for von Neumann algebras
\cite[p.~2]{AyupovKudaybergenov2014Matrix}.

The local version, where each single value is matched by a derivation, was
settled for von Neumann algebras earlier. Hadwin, Li, Li and Ma answer a
question of Kadison on the algebra of measurable operators (Theorem~\ref{thm:banach:local-measurable})
\cite[p.~1]{HadwinLiLiMa2013}.\begin{theorem}[{\cite[Theorem~1, p.~2]{HadwinLiLiMa2013}}]\label{thm:banach:local-measurable}
Let $M$ be a von Neumann algebra on a Hilbert space $H$. If the abelian
summand $M_1$ of $M$ is discrete, then every local derivation on the algebra
of closed densely defined operators affiliated with $M$ is a derivation.
\end{theorem}

The proof uses a separating family of central idempotents
\cite[Lemma~2, p.~2]{HadwinLiLiMa2013}. A local derivation that is a
derivation on each piece is a derivation on the whole
\cite[p.~2]{HadwinLiLiMa2013}. On a minimal projection $Q$ of a discrete
abelian summand the algebra $QM$ is one-dimensional, so the condition holds
there \cite[p.~3]{HadwinLiLiMa2013}.

The local derivations of an algebra need not be derivations one by one, but
as a set they again form a Lie algebra. Hao and Chen prove this for every
algebra over a field (Theorem~\ref{thm:banach:local-lie}) \cite[p.~1]{HaoChen2026Local}.

\begin{theorem}[{\cite[Theorem~2.8, p.~4]{HaoChen2026Local}}]\label{thm:banach:local-lie}
Let $A$ be an algebra over a field. Then the set $\operatorname{LocDer}(A)$
of local derivations of $A$ is a Lie algebra under the commutator bracket.
\end{theorem}

The proof passes through a reflexive hull, a single object that covers the
local derivations and the almost inner ones \cite[p.~1]{HaoChen2026Local}. The
same argument works for a family of multilinear operations, so the statement
is not tied to the associative product \cite[p.~5]{HaoChen2026Local}. The
authors also prove the structural conjecture of Ayupov, Elduque and
Kudaybergenov on local derivations \cite[p.~1]{HaoChen2026Local}.

A parallel line drops boundedness instead of linearity. For a type I von
Neumann algebra $M$ with center $Z$, let $LS(M)$ be the algebra of locally
measurable operators affiliated with $M$
\cite[p.~1]{AlbeverioAyupovKudaybergenov2007}. Albeverio, Ayupov and
Kudaybergenov ask for the derivations of $LS(M)$ that are linear over $Z$ (Theorem~\ref{thm:banach:zlinear-derivation})
\cite[p.~1]{AlbeverioAyupovKudaybergenov2007}.\begin{theorem}[{\cite[Theorem~3.9, p.~11]{AlbeverioAyupovKudaybergenov2007}}]\label{thm:banach:zlinear-derivation}
Let $M$ be a type I von Neumann algebra with center $Z$, and let $A$ be a
subalgebra of $LS(M)$ that contains $M$. Then every $Z$-linear derivation on
$A$ is spatial and implemented by an element of $LS(M)$.
\end{theorem}

The measurable operators and the totally measurable ones are such
subalgebras, so their $Z$-linear derivations are spatial too
\cite[Corollary~3.10, p.~11]{AlbeverioAyupovKudaybergenov2007}. When the
lattice of projections of $M$ is atomic, every derivation of the
$\tau$-measurable operators is spatial, hence continuous in the measure
topology \cite[Corollary~3.11, p.~11]{AlbeverioAyupovKudaybergenov2007}.

For the $\tau$-compact operators the same authors give the full description
\cite[p.~1]{AlbeverioAyupovKudaybergenovKalandarov2008}. Let $S_{0}(M,\tau)$
be that algebra (Theorem~\ref{thm:banach:tau-compact}).

\begin{theorem}[{\cite[Theorem~3.7, p.~10]{AlbeverioAyupovKudaybergenovKalandarov2008}}]\label{thm:banach:tau-compact}
Let $M$ be a type $\mathrm{I}_\infty$ von Neumann algebra with a faithful
normal semi-finite trace $\tau$. Then every derivation of $S_{0}(M,\tau)$ is
spatial and implemented by an element of $S(M,\tau)$.
\end{theorem}

Outside type $\mathrm{I}_\infty$ the centre of $S_{0}(M,\tau)$ can be
trivial, and the methods of the earlier papers do not apply
\cite[p.~10]{AlbeverioAyupovKudaybergenovKalandarov2008}.

The same circle of ideas returns for finite von Neumann algebras, where the
algebras of affiliated operators carry the quantum-mechanical picture that
opened this chapter \cite[p.~2087]{KadisonLiu2014Murray}. Kadison and Liu work with a finite von Neumann algebra
$R$ and its Murray--von Neumann algebra $A_f(R)$ of affiliated operators
\cite[p.~2089]{KadisonLiu2014Murray}. They study the derivations of
$A_f(R)$ that map $R$ into $R$ (Theorem~\ref{thm:banach:kadison-liu}) \cite[p.~2089]{KadisonLiu2014Murray}.

\begin{theorem}[{\cite[Theorem~5, p.~2090]{KadisonLiu2014Murray}}]\label{thm:banach:kadison-liu}
Let $\delta$ be an extended derivation of $A_f(R)$. Then there is an operator
$B$ in $R$ such that $\delta(A)=A\cdot B-B\cdot A$ for every $A$ in $A_f(R)$.
\end{theorem}

The proof restricts $\delta$ to $R$, applies
Theorem~\ref{thm:banach:sakaiinner} to obtain $B$, and then shows that
$\delta-\operatorname{ad}B$ kills all of $A_f(R)$
\cite[p.~2090]{KadisonLiu2014Murray}. For a factor of type $\mathrm{II}_1$ the
paper goes further: the only derivation of $A_f(M)$ with values in $M$ is zero
\cite[Theorem~12, p.~2091]{KadisonLiu2014Murray}. Kadison and Liu record that
the full statement, that every derivation of $A_f(R)$ is inner, is still open
\cite[p.~2090]{KadisonLiu2014Murray}.

Nonlinear maps on a factor von Neumann algebra give another extension of the
same theorem \cite[p.~1979]{YuZhang2012StarLie}. Yu and Zhang drop linearity and the involution is twisted in the
rule (Theorem~\ref{thm:banach:yu-zhang}) \cite[p.~1979]{YuZhang2012StarLie}.

\begin{theorem}[{\cite[Theorem~2.1, p.~1980]{YuZhang2012StarLie}}]\label{thm:banach:yu-zhang}
Let $M$ be a factor von Neumann algebra on a complex Hilbert space $H$ with
$\dim H\geq2$. Suppose that $\varphi\colon M\to M$ is a nonlinear $*$-Lie
derivation, that is,
$\varphi([A,B]^*)=[\varphi(A),B]^*+[A,\varphi(B)]^*$ for all $A,B$ in $M$,
where $[A,B]^*=AB-BA^*$. Then $\varphi$ is an additive $*$-derivation.
\end{theorem}

The proof splits each element over the two Peirce spaces of a projection and
checks the derivation rule on the four blocks, then uses
Theorem~\ref{thm:banach:sakaiinner} for the inner form
\cite[pp.~1990--1991]{YuZhang2012StarLie}.

\hbfigure{18-banach-operator/derivation-containment}{The derivation spaces of a $C^*$-algebra $A$ on a Hilbert space $H$. An inner derivation $x\mapsto ax-xa$ has $a$ in $A$, so it is spatial, and a spatial derivation is a derivation \cite[pp.~273--274]{Sakai1966}. Kadison proves that every derivation of a $C^*$-algebra on $H$ is spatial \cite[Theorem~4, p.~283]{Kadison1966}, and Sakai that every derivation of a von Neumann algebra is inner \cite[Theorem~1, p.~274]{Sakai1966}, so on a von Neumann algebra all three spaces agree \cite[Theorem~15, p.~290]{Kadison1966}.}{fig:banach:derivation-containment}

The norms of the inner derivations of a tensor product of $C^*$-algebras are related to the minimal prime ideal spaces of the factors~\cite{WlMinimalPrimeIdealSpacesAndNorm1996}.

Lazar, Tsui and Wright study when a derivation of a hereditary $C^*$-subalgebra $B$ of an approximately finite-dimensional $C^*$-algebra $A$ extends to a derivation of $A$, after a perturbation by a derivation generated by a multiplier of $B$~\cite{WlExtendingDerivationsFromHeredi1998}.

Hamhalter, Kudaybergenov, Peralta and Russo prove a Jordan version of Dorofeev's boundedness theorem for completely additive measures, and use it to show that every, not necessarily linear or continuous, 2-local triple derivation on a continuous $JBW^*$-triple is a triple derivation~\cite{WlBoundednessOfCompletelyAdditiv2015}.
